\scanpage{213}
\mspage{162}

où $i$ est un monomorphisme. On sait que le caractère cocartésien du carré se vérifie ``argument par argument'' par le caractère cocartésien de

\diagram{\begin{tikzcd}[ampersand replacement=\&]
\deleted{\unsure{F}}G(x)\arrow[r,hook,no head,"i(x)"]\arrow[d,"p(x)"'] \& F(x)\arrow[d,"q(x)"]\\
G'(x)\arrow[r,hook,no head,"i'(x)"] \& F'(x)
\end{tikzcd}}

dans Ens, $\forall x$ dans $X$. Or un \inserted{tel} carré cocartésien \unsure{sur|dans} $\unsure{\overline{\mathrm E}|\mathrm E}\mathrm{ns}$, avec $i(x)$ mono, est aussi cartésien. Cela implique que le carré $Q$ dans $X^\wedge$ est également cartésien.

Considérons \inserted{d'autre part} le foncteur canonique

\[
\varphi_X\colon F\longmapsto X/F\colon X^\wedge\longrightarrow\mathrm{Cat}
\]

Il vaut mieux, d'ailleurs, le regarder comme un foncteur

\[
\varphi_{\unsure{X|x}}\colon X^\wedge\longrightarrow\mathrm{Cat}_{/X},
\]

et je rappelle, pour mémoire, que ce foncteur est pleinement fidèle, et que son image essentielle est formée des catégories $X'$ sur $X$ telles que

\[
\unsure{\alpha'|X'}\mathrel{\unsure{-|\longrightarrow}}X
\]
% archive sheet number: 213
